\section{Finite memory, priors, and information}\label{l:model}

We specify the observation model once, before applying the different
localization estimates. A signal \(s\in S^{d-1}\) is fixed before
sampling. At index \(j\) the learner receives
\[
 x_j\sim N(0,I_d),\qquad y_j=\langle x_j,s\rangle,
\]
with independent rows \(x_j\). It cannot choose a row or revisit an
earlier observation. Its persistent state at each index lies in a
finite set of cardinality at most \(2^M\). Given that state and the
entire current pair \((x_j,y_j)\), a transition may use unrestricted
computation and fresh randomness. All information retained for later
indices must be represented by the new state.

The learner stops by a prescribed deterministic horizon \(T<\infty\),
possibly before taking a sample. Its unit-vector output uses only the
terminal state, the stopping index, and fresh randomness. The rules and
output laws may also depend on data-independent shared randomness
chosen before sampling, as well as on \(d\), the prescribed accuracy,
and the index; they cannot depend on the unknown signal. The shared
rule randomness and the initial finite state are jointly independent
of the signal and all sample rows. The initial state may depend on the
shared randomness, but conditional on that randomness it carries no
signal or sample information. In particular, the
last observation affects the output only through the terminal state.
This is a bit-state model, not a model with persistent real registers.

For a unit-vector estimate \(\widehat S\), angular success at
accuracy \(\epsilon\) is the event
\(\arccos\langle S,\widehat S\rangle\le\epsilon\).
It implies \(\|S-\widehat S\|\le\epsilon\), since the
chord length at angle \(\theta\) is \(2\sin(\theta/2)\le\theta\).
Each application states whether success is averaged under a specified
prior or required for every fixed signal.

As part of the randomized-learner model, the initial variables,
sequential finite choices, and output form a jointly measurable
experiment in the signal, the sample rows, and the randomness.
Consequently prior mixing and conditioning on the shared randomness
refer to measurable joint laws. The next regularity allowance concerns
the coordinate kernels of that experiment at a fixed shared seed.

For each fixed realization of the shared rule randomness, the
transition and stopping probabilities may be measurable in the
completion of the Borel sigma field under
\[
 \lambda(dx,dy)=\gamma_d(dx)\,dy,
\]
where \(\gamma_d\) is standard Gaussian probability on \(\R^d\).
The terminal output laws are Borel probability measures on the sphere.
The next lemma explains why reference experiments may nevertheless
evaluate state rules at every matrix and label.

\begin{lemma}[Borel rules under an absolutely continuous prior]
\label{l:borel}
Let \(P\) be a probability measure on \(S^{d-1}\) such that, for
independent \(S\sim P\) and \(x\sim\gamma_d\), the pair
\((x,\langle x,S\rangle)\) has law absolutely continuous with respect
to \(\lambda\). Fix a realization of the shared rule randomness and
a finite horizon. Every completed-measurable finite-state learner has
everywhere-defined Borel transition and stopping kernels with the
same state sets and output laws and the same \(P\)-prior joint law of
the signal, state path, stopping index, and output.
\end{lemma}

\begin{proof}
At a fixed state and index, combine continuation and stopping choices
into one finite set of tagged destinations. Choose Borel versions of
their completed-measurable coordinate probabilities. Such versions
exist: measurable simple approximations may have each of their
countably many level sets replaced by a Borel set modulo a null set.
Outside a Borel \(\lambda\)-null set the versions agree with the
original coordinates and form a probability vector. On the exceptional
set use a fixed point mass. This gives nonnegative Borel coordinates
whose sum is one everywhere.

There are finitely many states and indices, so a single Borel
\(\lambda\)-null set contains every modification. Pregenerate all
sample pairs through the horizon, including pairs ignored after
stopping. Each of their unconditional marginals avoids this null set
by the hypothesis on \(P\). A finite union over indices therefore
shows that no modified pair occurs almost surely. This argument uses
the unconditional pair law; it does not require a positive density
after conditioning on a state selected from past samples.

Couple the old and new finite kernels by the same fresh uniforms.
Induction gives identical state paths and stopping decisions almost
surely. At an identical terminal state and index the output laws are
unchanged and can be coupled identically. This proves the joint-law
claim.
\end{proof}

\begin{lemma}[Sphere and regular cube-chart priors]\label{l:prior-ac}
The hypothesis of Lemma~\ref{l:borel} holds for uniform probability
\(\sigma\) on \(S^{d-1}\), \(d\ge2\). It also holds for
\(S=\phi(Z)\) when \(Z\) has a density relative to Lebesgue measure
on a bounded \(q\)-dimensional cube, \(q\ge1\), and \(\phi\) is
\(C^1\) on a neighborhood of its interior with
\(\operatorname{rank}D\phi(z)\ge1\) for almost every \(z\).
\end{lemma}

\begin{proof}
For fixed \(x\ne0\), rotational invariance makes
\(\langle x,S\rangle/\|x\|\) the first coordinate of a uniform
spherical point. Its density on \((-1,1)\) is proportional to
\((1-t^2)^{(d-3)/2}\); this also is integrable for \(d=2\).
For example the formula follows by writing the point as a normalized
standard Gaussian and separating its first coordinate from the
remaining squared length. Thus the conditional label law has a
Lebesgue density for every nonzero \(x\). Fubini, discarding \(x=0\),
proves the first assertion.

For a chart put \(f_x(z)=\langle x,\phi(z)\rangle\). At every
regular \(z\), the set of \(x\) with
\(\nabla f_x(z)=D\phi(z)^{\mathsf T}x=0\) is a proper linear
subspace and has Gaussian measure zero. Fubini shows that for
Gaussian-almost every \(x\), the critical set of \(f_x\) has
Lebesgue measure zero in the cube. Around every other interior point,
choose a coordinate in which a partial derivative is nonzero. The
inverse function theorem makes \(f_x\), together with the remaining
coordinates, a \(C^1\) local coordinate system. In each such
neighborhood the preimage of a Lebesgue-null subset of \(\R\) has
zero \(q\)-dimensional volume, by Fubini and change of variables.
A countable cover of the regular set proves absolute continuity of
the label pushforward. The critical set and cube boundary are null.
Integrating over \(x\) proves the second assertion.
\end{proof}

When a proof fixes randomness under a chosen prior, the order matters.
If the original success is at least \(p\) and \(q<p\), first fix the
shared rule randomness with prior success greater than \(q\).
Apply Lemma~\ref{l:borel} for that fixed rule and that prior. Then,
the conditional initial-state law remains independent of the signal
and rows. If a deterministic program is useful, sample that initial
state, sample the finite choices by uniforms indexed by state and
index, and sample the output law for each terminal state and index.
These variables are independent of the inputs. Averaging fixes them
with prior success at least \(q\).

Alternatively, when the same numerical success bound is proved at
each fixed shared seed, apply the Borel replacement separately there
and transfer the bound back to the original conditional prior law.
The original jointly measurable experiment makes those conditional
success probabilities measurable, so their bound may be integrated
over the seed. Neither argument requires a jointly chosen family of
Borel replacements. A replacement for one prior preserves that prior
experiment, rather than every individual signal.

\paragraph{Blocks and the stopping index.}
Starting from a known state, composition of the everywhere Borel
kernels gives a Borel probability kernel on the complete data of a
block. A block of \(b\) rows can report its continuing state, or its
terminal state and a stopping offset in \(\{0,\ldots,b\}\), or a
dummy value if it was already inactive. Thus
\[
 N_b\le (b+2)2^M+1
\]
is a valid local alphabet bound. A proof may retain a larger displayed
bound from its own convention. The sequence of block indices and local
offsets determines the global stopping index. Pregenerating unused
suffixes keeps every new block independent of the past. Auxiliary
locations or projections introduced below are part of the analysis;
they are not supplied to the learner.

\paragraph{Information conventions.}
All logarithms are natural unless a dyadic scale is explicitly written
with \(\log_2\). For probabilities \(P\ll Q\) on a standard Borel
space,
\[
 D(P\|Q)=\int\log(dP/dQ)\,dP,
\]
with value \(+\infty\) without absolute continuity. For a discrete
variable \(W\), \(H(W)=-\sum_w P(W=w)\log P(W=w)\), where
\(0\log0=0\). Mutual information is
\(I(U;W)=D(P_{U,W}\|P_U\otimes P_W)\); conditional quantities average
the corresponding conditional relative entropies. We use the chain
rule and data processing on standard Borel spaces in the
conditional-relative-entropy formulation of
\cite[Lemma~3.6 and Sections~3.3--3.4]{Austin2020}.

Two elementary consequences will be used repeatedly. If an event has
probability \(b\) under \(P\) and \(b_0\) under \(Q\), then
\begin{equation}\label{l:binary}
 D(P\|Q)\ge b\log(1/b_0)-\log2.
\end{equation}
Indeed data processing gives binary relative entropy, and the sum of
\(b\log b+(1-b)\log(1-b)\) is at least \(-\log2\), while the
remaining term \(-(1-b)\log(1-b_0)\) is nonnegative. If \(P\ll Q\),
\(a>0\), and \(F>0\), then
\begin{equation}\label{l:tilt}
 \E_P\log F\le
 \frac{D(P\|Q)+\log\E_QF^a}{a}
\end{equation}
whenever the left side is defined and the right side is finite.
For a bounded logarithm, normalize \(F^aQ\) and use nonnegativity of
relative entropy. Apply this argument to clipped logarithms and use
monotone truncation of the positive part and integrability of the
negative part for the displayed form. This is the entropy
variational inequality; the bounded-measurable probability version
is recorded in \cite[Equations~(1.1)--(1.2)]{DupuisMao2022}.

Relative entropy to a finite nonzero measure \(\eta\) means
\(\int\log(dP/d\eta)\,dP\). If \(c=\eta(\Omega)\), it equals
\(D(P\|\eta/c)-\log c\). Therefore it contracts under a measurable
map: normalization adds the same constant before and after the map.
This observation is used only where the comparison measure and all
relevant logarithmic integrals are specified.

\begin{lemma}[Spherical caps and a full-data endpoint]\label{l:caps-endpoint}
Write \(n=d-1\). Uniform spherical probability satisfies
\[
 \sigma(B(u,r))\le r^n\quad(u\in S^{d-1},\ r>0),
 \qquad
 \sigma(B(z,r))\le(2r)^n\quad(z\in\R^d,\ r>0).
\]
If \(S\sim\sigma\), \(0\le T<d\) independent Gaussian rows and all
their exact labels are revealed, and an estimate uses these data and
independent randomness, then for every \(0<\epsilon<r_0<1\),
\begin{equation}\label{l:raw-success}
 \Prob\{\|\widehat S-S\|\le\epsilon\}
 \le \frac12+\frac12\,
 \Prob\!\left\{\|P_{\ker X}S\|\le r_0\right\}
 \le \frac12+\frac{T}{2d(1-r_0^2)}.
\end{equation}
\end{lemma}

\begin{proof}
For \(0<r\le1\), a chordal cap about \(u\) has angular radius
\(\theta\le\pi/3\), \(\sin\theta\le r\), and
\(\cos\theta\ge1/2\). Its normalized area is
\[
 \frac{\int_0^\theta\sin^{n-1}v\,dv}
      {\int_0^\pi\sin^{n-1}v\,dv}.
\]
The numerator is at most \(2r^n/n\) after substituting
\(w=\sin v\); the denominator is at least \(2/n\), since
\(\int_0^{\pi/2}\sin^{n-1}v\,dv
\ge\int_0^{\pi/2}\sin^{n-1}v\cos v\,dv=1/n\).
For \(r\ge1\) the same bound is trivial. An ambient ball meeting the
sphere is contained, on the sphere, in a radius-\(2r\) ball about
any point of its intersection, proving the second bound.

The row span of \(X\) has dimension \(T\) almost surely. Its labels
specify the row-space projection of \(S\). Conditional on the rows
and labels, the remaining component is uniform on its sphere in
\(\ker X\), with radius \(R=\|P_{\ker X}S\|\). To see this, write
\(S\) as a normalized standard Gaussian and separate its independent
row-space and nullspace components; the direction of the latter is
independent of both lengths and of the former component.

When \(R>r_0>\epsilon\), the set of residual directions that place
the signal within \(\epsilon\) of a fixed estimate cannot contain
both members of any antipodal pair. Its conditional mass is at most
\(1/2\), including after averaging the estimate's independent
randomness. Finally,
\(\E\|P_{\operatorname{row}X}S\|^2=T/d\) by rotational
invariance. Markov's inequality applied to
\(\|P_{\operatorname{row}X}S\|^2=1-R^2\) proves
\eqref{l:raw-success}.
\end{proof}

The constants used later follow directly from this proof. For success
at least \(2/3\) and \(\epsilon\le1/10\), the same calculation with
the strict event \(R<1/2\) gives \(T>d/4\): if \(T\le d/4\),
then \(\Prob(R<1/2)<1/3\) and success is less than \(2/3\).
The weaker \(T>d/8\) follows by taking \(r_0=1/\sqrt2\), which
gives success at most \(5/8\) when \(T\le d/8\).
For a fixed program with success at least \(3/5\), take the event
\(R\le\epsilon\) directly. Its probability is at most
\(T/[d(1-\epsilon^2)]\), so \(T\le d/20\) would give success
strictly below \(3/5\). These statements grant the entire stream
and therefore include early stopping.
